\documentclass[12pt]{article}
%\usepackage{NJDnatbib}
\setlength{\textwidth}{6.3in}
\setlength{\textheight}{8.7in}
\setlength{\topmargin}{0pt}
\setlength{\headsep}{0pt}
\setlength{\headheight}{0pt}
\setlength{\oddsidemargin}{0pt}
\setlength{\evensidemargin}{0pt}

%\usepackage[utf8]{inputenc}

\usepackage{amssymb}

\usepackage{amsmath,amsthm,amssymb}

\usepackage[T1]{fontenc}
%\usepackage{graphicx}

\usepackage[colorlinks=true,citecolor=black,linkcolor=black,urlcolor=blue]{hyperref}
%\newcommand{\arXiv}[1]{\href{http://arxiv.org/abs/#1}{\texttt{arXiv:#1}}}


\theoremstyle{plain}
\newtheorem{thm}{Theorem}

\newtheorem{lem}{Lemma}[section]
\newtheorem{prop}[lem]{Proposition}
\newtheorem{corollary}[lem]{Corollary}
\newtheorem{claim}[lem]{Claim}
\newtheorem{fact}[lem]{Fact}

\theoremstyle{definition}
\newtheorem{definition}[lem]{Definition}
\newtheorem{observation}[lem]{Observation}
\newtheorem{remark}[lem]{Remark}

\newbox\PPbox % symbol for ++
\setbox\PPbox=\hbox{\kern.5pt\raise1pt\hbox{\small+\kern-1pt+}\kern.5pt}
\def\PP{\copy\PPbox}
\newcommand\Cee{\textsc{C}}
\newcommand\Cpp{\Cee\PP}
\newcommand\foreign[1]{\textsl{#1}}

\newcommand\std[1]{\texttt{std::#1}}
\def\frwl.{\std{forward\_list}}
\def\uptr.{\std{unique\_ptr}}
\def\insa.{\textit{insert\_\!after}}
\def\ersa.{\textit{erase\_\!after}}
\def\spla.{\textit{splice\_\!after}}
\def\slnd.{\textbf{\textit{sl\_\!node}}}
\def\lnkt.{\textbf{\textit{link\_\!type}}}
\def\lnlc.{\textit{link\_\!loc}}
\def\slls.{\textbf{\textit{sl\_\!list}}}
\def\smpl.{\textbf{\textit{simple\_\!list}}}
\def\mspl.{\textbf{\textit{mirrored\_\!simple\_\!list}}}
\def\msll.{\textbf{\textit{mirrored\_\!sl\_\!list}}}
\def\atend.{\textit{at\_\!end}}
\def\queue.{\textbf{\textit{queue}}}
\def\stack.{\textbf{\textit{stack}}}

\def\ang#1{\langle#1\rangle}

\title{Singly linked lists as a first class C++ container type}
\author{Marc A. A. van Leeuwen\\
\small Laboratoire de Math\'ematiques et Applications\\[-0.8ex]
\small Universit\'e~de~Poitiers\\[-0.8ex]
\small Poitiers, France\\
\small\tt Marc.van-Leeuwen@math.univ-poitiers.fr
}

%\date{\dateline{13 Oct 2017}\\ \small Classifications:}

\begin{document}

\maketitle

\begin{abstract}
The \Cpp\ container library provides various general container class templates,
which can be classified into containers with constant time random access
(\texttt{array}, \texttt{vector}, and \texttt{deque}) and those using linked
structures (all others). Nearly all linked structure containers require at least
two links per node; only \texttt{forward\_list} (which was added in the \Cpp11
revision of the language) provides singly linked lists. The specification of
\texttt{forward\_list} is awkward however, with iterators that are to
be used in a different manner from those in any other container class, and
notably with a distinction on whether they are used just for accessing nodes or
also for changing the structure of the list itself (insertion/deletion/splicing
of nodes).

We present a container class template for singly linked lists with regular
forward iterators. These can be used like those in any container class, notably
for purposes of node insertion and deletion. They are implemented using a well
known technique, namely as pointers to link fields. Appropriate class
templates have been developed in the context of the Atlas of Lie Groups and
Representations project, and used in many parts of its software program. We here
present our experience from this usage, which is largely satisfactory. Singly
linked list provide a simple, lightweight and flexible programming tool
in programming, and provide a more lightweight implementation for stacks an
queues than the \Cpp~library currently provides. A useful aspect is that there
are some variations of the data structure with easy conversions between them;
this allows for instance using lists even in places where using other container
types would be difficult. They seem a useful addition to the programmer's
toolbox.

\end{abstract}
\maketitle


\section{Introduction}

In the practice of software development in \Cpp, arguably the most important
difference with respect to using the \Cee\ programming language is having
available the class templates of the standard container library. In our ``Atlas
of Lie groups and Representations'' software, a project that has been in
development since 2004 and the context in which the ideas and techniques
discussed in this paper were developed, we choose in almost all cases to use the
class template \std{vector} rather than native \Cpp\ arrays. Since our software
operates on mathematical structures specified by the user at run time, it is
almost never possible for the size of composite values to be usefully fixed or
bounded at compile time, which reason alone justifies our choice in most cases.
Without container class templates, one would need to explicitly implement the
dynamic memory management provided by \std{vector}, either separately every time
a new type is involved, or (if using generic pointers) introducing type
insecurity that needs shielding off; neither of these are nearly as convenient
to the programmer as using the vector class template.

With speed of random access almost on a par with that of native arrays, and with
the added possibility to grow and shrink slots in amortised constant time, the
vector class template is a good choice in many cases where one needs to store a
sequence of values whose size is unknown at compile time. But there are of
course purposes for which this type of container is less suited, and the
\Cpp\ standard library provides several other container class templates that
address other needs: \texttt{array}, \texttt{deque}, \texttt{forward\_list},
\texttt{list}, \texttt{set}, \texttt{map}, \texttt{multiset}, \texttt{multimap},
\texttt{unordered\_set}, \texttt{unordered\_map}, \texttt{unordered\_multiset},
\texttt{unordered\_multimap} (all inside the \texttt{std} namespace).
The first two of these provide constant time random access, and can be
considered to be variations of native arrays respectively of \texttt{vector};
the remaining ones all to some extent use linked data structures (although for
the ``unordered'' versions the primary access is through hashing). The templates
\texttt{set}, \texttt{map} and their variations define associative containers,
which are primarily accessed by looking up \emph{key} values; in this article
our main interest will be in \emph{sequence} containers, where access is by
position within the (linear) container.

In linked data structures most items cannot be accessed directly, but must be
reached via links through a sequence of other items; this means there is no
random access in constant time. Linked sequential structures are most suited for
applications where access is naturally sequential (either one directional or
going back and forth), especially when the size of the sequence varies. Such
applications can be quite important though: for instance, \textsc{Lisp} and many
functional programming languages are largely focused on handling linked
structures, and in the \TeX\ typesetting engine most work involves handling
linked lists of tokens, boxes, etc. In algebraic the core of our `Atlas'
software, occasions where using linked lists are a natural choice are not as
frequent as those where it is natural to prefer structures based on
\std{vector}, but they certainly do occur. To name a few: some algorithms like
orbit generation or similar closure processes need a stack or queue of ``yet to
consider'' items; we often need a temporary container to pass a collection of
values computed by one function to be processed by another. The software also
supports a full user programming language, and its various parts (parser, type
checker/compiler, and runtime execution system) very frequently employ linked
structures of various kinds.

Nonetheless, the linked list container class templates provided by the \Cpp.
standard library appear to be not nearly as useful as \std{vector} is. In our
project, we actually \emph{never} opted to use \std{list}. Apart from
\std{vector} we made occasional use of \std{set} and \std{map}, and sometimes
defined a recursive type to implement a linked structure ``manually'', for
instance for representing a parse tree. Reasons for avoiding most linked
container types are subjective, but an important sticking point was that
\std{list} provides \emph{doubly} linked lists, which invariably was more than
we needed and therefore sub-optimal in space and time. As an indication that we
are not alone in our reluctance to use \std{list}, note that in the standard
container library itself the container adaptors \std{queue} and \std{stack},
which can perfectly well be implemented using \std{list}, prefer to be based by
default on \std{deque} instead.

The \frwl. class template, which was introduced in the \Cpp11 revision of the
language definition, does provide a singly linked list container type, and shows
that the language designers considered that just providing doubly linked lists
is not entirely satisfactory. Indeed, their motivation is clear in the report:
``It is intended that \texttt{forward\_list} have zero space or time overhead
relative to a hand-written C-style singly linked list. Features that would
conflict with that goal have been omitted.''. And the possible utility of such a
container type is witnessed by the notion of ``forward iterator'', a weak form
of the iterator abstraction that omits any functionality that singly linked
lists cannot efficiently provide (notably decrementing); this notion was already
present in the \Cpp03 specification even though, prior to the advent of \frwl.,
none of the container types in the standard actually provided iterators of
(just) this category.

However, the specification of \frwl. is manifestly awkward, with member
functions such as \textit{before\_begin}, \insa. and \ersa. not present for any
other container type; clearly \texttt{forward\_list} requires handling iterators
in an unusual manner. Again a motivation is provided in the report, namely
``Modifying any list requires access to the element preceding the first element
of interest, but in a \texttt{forward\_list} there is no constant-time way to
access a preceding element''; this refers to the impossibility to (efficiently)
perform the insertion of a node before a given node, or the deletion of the node
itself, given only a pointer to a node in a singly linked list.

But that indicates a narrow focus on one possible implementation of iterators
into a list, namely by pointers to nodes, while the iterator abstraction in no
way forces such an implementation. Indeed there is a well known alternative
method for list traversal that does not have this limitation, which instead uses
pointer to link fields (either within a node or at the head of the list), which
fields contain pointers to nodes; in other words, one adds a level of
indirection. This observation inspired us to define the container class template
discussed in this paper, for our own practical use and also as an experiment to
see if there were any unforeseen problems preventing the use of a singly list
container type in the same way as any other container type. No major
difficulties were found. As our work did not involve any truly novel ideas, we
remain somewhat surprised at not having found any indication that the
\Cpp\ community ever attempted a similar avenue.

In this paper we shall first give some technical background to singly linked
lists and iterators for them, and then proceed to describe our implementation of
a container class template using singly linked lists. We discuss how these lists
are used within our software package, and some adaptations and variation that
were inspired by practical questions arising while using them. Our lists can be
used as a fully capable container type in conjunction with the standard library,
but occasionally a somewhat less abstract approach proved to be useful. Our new
container blends in perfectly with using the standard library.

We illustrate these matters by concrete examples taken from the `Atlas'
software. Our implementation is in the public domain, licensed under the GNU
Public Licence, and available from our source code repository \cite{Atlas
  source}.

The \Cpp\ language is constantly evolving, but in this paper we shall take the
\Cpp11 revision as our reference point. This corresponds roughly to the current
state of our project (although we has since moved to requiring \Cpp14 support),
being a compromise between adopting useful innovations and not restricting our
software to users that have the most recent compiler versions installed.

\section{Design}
\subsection{Iterators in \Cpp.}

In \Cpp, the notion of iterator derives from and generalises the way in which
(already in~\Cee) pointer arithmetic serves a role in accessing native arrays:
the array provides a pointer, which is then modified using pointer arithmetic
and finally dereferenced do give access to an array element. This is a low-level
mechanism (and it exposes the user to the risk of out-of-bounds access), but its
generalisation to iterators associated to containers serves as a powerful
abstraction, allowing different container types, as well as user defined types
that adhere to the conventions, to be treated uniformly for many purposes. The
idea is that whatever the nature of the container, it will (amongst others)
provide access to its elements through values of an intermediary ``iterator''
type, which at a minimum permits to visit the elements of the container
sequentially in a fixed order. (This limited use of iterators generalises just
the ``increment'' form of pointer arithmetic.)

As a bonus, any pair$(p,q)$ of iterator values obtained from a same container
value will delimit a range of its elements, namely those accessed using the
iterator~$p$ by incrementing it a number of times without becoming equal to~$q$,
and then dereferencing it. Such a range can for many purposes be used instead
of the full contents of the container. In this role, iterator values serve as
\emph{fence posts between elements} in the container: the indicated range
consists of all items between the fence posts $p$ and $q$. Dereferencing an
iterator accesses the item after the corresponding fence post: dereferencing $p$
gives the first element of the range, while dereferencing $q$ gives the first
element \emph{after} the range. The fence post indicating the very end of the
container is called the ``past the end'' iterator, and cannot be validly
dereferenced; its existence is mandated uniquely to serve as ending fence post.

Thus a container with $n$ elements has $n+1$ valid iterator values associated to
it, one before each item in the container, and an extra one at the end. The
usual terminology is that an iterator ``points to'' or (in the standard)
``refers to'' the element that would be obtained by dereferencing it, and the
syntax for the access of elements through iterators is the same as the syntax
for dereferencing pointers. However we find this terminology too suggestive of a
particular implementation of iterators; to stress the abstract relation, we
shall instead say that an iterator ``gives access'' to the element after the
fence post that the iterator corresponds to, or to no element at all in case of
a past-the-end iterator.

This ``in between elements'' nature of iterators can also be seen in the notion
of reverse iterators. These can be defined only for certain types of containers,
but when they can, each iterator value~$p$ has a corresponding reverse iterator
value that we shall denote here by $\tilde p$ (the actual syntax for them is of
no importance here), which contains the same information as~$p$ but has the
opposite orientation: incrementing $\tilde p$ corresponds to decrementing~$p$.
This for uses that require just a pair of iterators to describe a sequence of
values, the values in the range determined by~$(p,q)$ can be interpreted as a
sequence in the reverse order by specifying the range by $(\tilde q,\tilde p)$
instead. In particular $\tilde p$ gives access to the item \emph{before} the one
$p$ gives access to, if there is such an item; otherwise $p$ is at the lower
extremity of the container, in which case $\tilde p$ is the past-the-end reverse
iterator.


The fence post interpretation also applies in a third kind of iterator usage,
only possible with containers that can change size, namely to locate a position
for inserting one or more items into the container. Insertion can be done either
between two successive items, or at one of the extremities of the sequence, in
other words at every valid iterator value. (Similarly removal of a sequence of
items can be indicated by a pair of iterators, the sequence being the range
determined by the pair as indicated above.) As above, the insertion point
indicated by an iterator is situated just before the item it gives access to.

Any operation that changes the size of a container also changes the number of
valid iterator values for it in the obvious way. Whether all or any of the
previously valid iterator values remain valid, and if so which position in the
new container they correspond to depends very much on the container type and
details of the situation. The case of a vector requiring reallocation to fresh
memory is one in which all iterators are invalidated and a whole new slate of
iterator values comes into existence. The \Cpp. standard specifies in a case by
case fashion which iterators can be invalidated by which operations for a given
container class, which mostly reflects what happens for the envisioned
implementation of iterators. It seems fair to say however that there is no
general expectation users should have in this respect, that hold independently
of implementation considerations; it is not even true that an iterator that
remains valid necessarily gives access to the same item as it did before the
operation. As a general rule, insertion or removal operations in linked data
structures tend to only invalidate iterator values directly involved
in the operation.

It occurs frequently that one container class allows for iterators of more than
one type; this is unproblematic as long as different types are not mixed (for
instance in delimiting a range). For instance arrays and vectors have their
abstract iterator types, but also allow using raw pointers to elements as
iterators; we also saw the coexistence of forward and reverse iterators in some
cases. In such cases, what was said about the number of valid iterator values
applies to iterators of any one type, and between two different iterator types
for the same container, the sets of these values are in an obvious bijection.

\subsection{Singly linked lists, and \frwl.}

We now turn our attention specifically to singly linked lists, and iterators for
them. Such lists are the simplest possible linked data structures; as such they
are formed by a collection of separate ``nodes'' that may reside in unrelated
places in memory, but which are logically linked together using pointers. Being
singly linked means each node contains a unique field (the ``link'' field, which
we shall call \textit{next}) containing a pointer to another node, which is its
successor in the list; the final node in the list has a null pointer in its link
field. In order to cater for the possibility of lists being empty, the initial
data that giving access to a list is usually a pointer rather than a node; this
``head'' pointer is the initial link for the list; it either refers to the
initial node in the list, or else holds a null pointer to indicate that the list
is empty.

With such a data structure, it takes a lot more work to find the predecessor of
a given node than to find its successor, so it is reasonable that an iterator
for singly linked lists should implement an increment operation but no decrement
operation (and no reverse iterators): such iterators are called forward
iterators. Forward iteration over a linked list whose structure remains fixed
can be done in an obvious way: a pointer is initialised with the head pointer of
the list; as long as this pointer is not null, it gives access to a node, and
the pointer can be incremented by assigning to it the \textit{next} field of
that node. Pairs of such pointers determine a range in the manner described
above; iteration over such a range is done just like iteration over a complete
list, except that comparisons against the null pointer are replaced by
comparisons against the ending pointer for the range.

There is however a difficulty with the third mentioned purpose for iterators, in
the context of insertion and removal of (ranges of) nodes. Given a pointer
value~$p$ occurring in a singly linked list, insertion of one or more nodes in
the list at the position indicated by~$p$ requires replacing the copy of~$p$
that occurs \emph{as a link in the list or as its head pointer} by a pointer to
a fresh node, whose \textit{next} field is initialised by another copy of~$p$.
Knowing the value of$~p$ does not give easy access to the location where this
new pointer must be stored; in case it is a \textit{next} field, that field
belongs to the predecessor of the node that $p$ points to, and that predecessor
can only be reached by iteration starting from the list head. Similarly, when a
node (or range of nodes) has to be deleted, a pointer~$p$ to that node
(respectively to the first node of that range) does not provide access to the
link that needs to be changed; again it is the link that holds the value~$p$,
either in the \textit{next} field of the previous node or the list head itself,
that needs change.

These issues are of course well known, and to a large extent independent of
using any notion of iterators. One way to deal with them is to use, when
handling a list with the intention to insert or remove nodes, a pointer to the
node \emph{before} the place where insertion or deletion is to take
place\footnote
%
 {Cases where all insertions and deletions are done at the front of the list
  are an exception to this; here one can directly operate on the list head.
  Thus a list can implement a stack without needing any pointer variables other
  than the list head. Our discussion focuses on more general situations.}.
%
Something then needs to be done to handle the case that the insertion or
deletion happens to be at the front of the list, for which no such node exists.
One could use the null pointer to signal this special case, as it cannot occur
otherwise. But another solution is to use the convention that the head node of
the list is stored as the \textit{next} field of a fake node, the rest of which
is unused and which is not counted as element of the list, but which can be the
target of a pointer. The latter approach has the advantage that no test for
insertion or deletion at the front of the list is needed, as it becomes a
regular case. As an example the \TeX\ engine, which was originally written in
Pascal and uses a large array global array~\textit{mem} as a pool for storing
list nodes, adopts this approach. In lieu of pointers it uses indices into a
global array~\textit{mem}, so its fake head nodes cannot reside in local
variables; they are either statically reserved locations in \textit{mem} like
$\textit{mem}[\textit{temp\_head}\,]$ (where \textit{temp\_head} is a constant),
or locations dynamically allocated by functions such as \textit{scan\_toks} that
build up new lists from scratch (see \cite{TeX the program} for details). Thus
each list occupies one more slot in~\textit{mem} than it has elements.

These ideas are echoed in the specification of the class template \frwl.. As can
be expected it has forward iterators, which can be used for front-to-back
traversal, or to designate ranges in the usual fashion for algorithm doing
things like copying or linear searching. But for insertion and removal, as well
as for the linked-list specific operation of splicing (which moves around ranges
between lists by adjusting some links), it imposes unusual conventions that are
clearly motivated by implementation considerations like those given above.
Notably, there are methods \insa. and \ersa. that require an iterator giving
access to the node \emph{before} the point of insertion or removal. This
iterator passed to these methods should therefore never be the ending iterator
of the list; however, it may be a special iterator value returned by the method
\textit{before\_begin} when insertion or deletion at the front of the
list is needed. This iterator does not access any node, and incrementing it will
result in the ordinary starting iterator for the list (the one returned
by~\textit{begin})\footnote
%
{Actually the \Cpp\ standard contradicts itself when it requires that
 \texttt{before\_begin} return ``a non-dereferenceable iterator that, when
 incremented, is equal to the iterator returned by \texttt{begin()}'' while
 requiring as precondition for the meaning of \texttt{\PP r} for an iterator
 expression~\texttt r that ``\texttt r is dereferenceable''.
}.
%
We see that a \frwl. containing~$n$ items has $n+2$ valid iterator values (of
which those at both extremities are not dereferenceable), contrary to the
situation for all other container types.
For \ersa., the end of the range to erase by contrast has the usual
interpretation (if it accesses a node, that node is just beyond the range); as a
consequence, the distance between the iterators used to specify this range is
one more than the number of nodes to erase.

Splicing is somewhat like insertion of a range of nodes from another list (or
the same list, as long as the range does not contain the point of insertion),
but no new nodes are created as the range is just un-linked from its original
position and linked into the destination position. To specify the general form
of splicing involves three iterators: one in the destination list for the
insertion position and two in the source list to specify a range. Of these the
insertion position and the first iterator for the range follow the ``one place
before'' convention, presumably in order to be able to do the link adjustment
efficiently, but the end of the range follows the ordinary convention (so that
again the number of nodes moved is one less than the distance between the
iterators for the range). The latter is curious, as the link across the end of
the range also needs assigning to, and if an iterator only gives access to the
link \emph{after} the node it points to, finding this link still requires
traversing the range; indeed the specified complexity is the size of the range,
and not constant time as one would expect had the ``one place before''
convention been used for all three iterators rather than just two of them.

All in all of iterators for \frwl. are specified in a way that, while
understandable, is not very elegant and differs from that for iterators for
other containers. Working with them requires some special care: one should avoid
accessing the fake node before the list (which would give undefined behaviour;
anyway the implementation probably does not even allocate any memory for it
other than for its \textit{next} field), and often iterating over a list with
the goal of inserting or deleting nodes requires a \emph{pair} of iterators: an
ordinary iterator to access the node, and another one lagging one step behind,
in order to have its value available when the place for the operation has been
found. All in all two different kind of iterators need to be distinguished
according to their intended use, but the type system cannot help enforce using
them appropriately.

\subsection{Pointers to pointers, for simpler iterators}

There is however another way around the difficulty that holding a pointer to a
node in a singly linked list does not allow changing the link to that node
inside the list structure. Rather than working with (a local copy of) such a
pointer, one can record the \emph{address} of the link in the list that holds
this pointer value, either the \textit{next} field of the previous node or the
head pointer for the list in case the node is the first one; the type of such a
variable is pointer to pointer to node. Access to the current node requires more
indirection with this approach (getting the pointer at that address), but it is
now also possible to assign to the pointer at that address, which changes the
list at the link towards the current node; this is exactly the functionality
that was missing above. Assuming that pointers are typed (their dereferencing
gives a value of a well determined type rather than a access to raw memory; we
assume this throughout), this technique requires two characteristics of the
pointer model: a same pointer can be made to point to a variables on the runtime
stack as well as to ones allocated on the heap, and it can point to a part of a
structure as well as to an isolated variable. Indeed list heads will be isolated
variables on the runtime stack, while other link fields are part of a node
allocated on the heap. The model of Pascal does not have these properties, since
it allows pointers only to point entire structures generated on the heap; by
contrast the model of \Cee\ (and \Cpp) has these characteristics. To our
knowledge the first high level programming language that allows this technique
is \textsc{Algol~68}, where it was known as the ``triple ref
technique''\footnote {That name derives from the fact that the mode of the
variables involved is \textbf{ref ref ref node}, where each reference level
corresponds to a level of mutability: the final one for modifying the node
contents, the preceding one for modifying links, and the initial one for
modifying the variable itself to traverse the set of link fields. In \Cee\ the
latter is implicit, so there the type of the variable has just two pointer
levels. } (see~\cite{triple ref} for a witness to this fact, and a popular
explanation, in contemporary social media).

In \Cpp\ it should therefore be possible to implement iterators for simply
linked lists in such a way that obviates the ``one step behind'' convention.
With an iterator holding the address of a link field, the dereference and
increment operations can start by fetching the link itself, and then get the
node contents parts or the address of its link part, respectively. For the
structure changing operations like node insertion, one instead uses the address
held to modify the link stored there. All allowed operations work uniformly,
whether the iterator is the initial, the final, or an intermediary one in the
list (of course dereference and increment are forbidden for the final iterator:
the link it refers to holds a null pointer).

We shall report below on details of the implementation in our software, and on
our experiences. In short, we shall see that iterators are used just like those
for standard library containers, that ordinary \textit{begin}, \textit{insert}
and \textit{erase} are used instead of \textit{before\_begin}, \insa. and
\ersa., and that instead of \spla. we have \textit{splice} methods that give
the ordinary interpretation to all their (up to~$3$) iterator arguments.

Questions of what happens with iterator values when structure changing operation
are applied to a list have more subtle answers, which do reveal some information
about the mechanism by which iterators are implemented: these answers are not
always the same as they are for iterators of say \frwl. or \std{list}. Our
iterators are attached to the node before rather than after the fence post they
represent (or to the list head itself if there is no such node), which differs
from the node that they give access to. This means that when the structure of
the list is modified, the iterator remains valid if and only if the node it was
attached to remains in existence. After the modification the iterator will give
access to whatever node now follows that node (if any); this is not necessarily
the one it gave access to before the modification. The iterator value attached
to the list head remains valid as long as the list object itself exists, but
during its lifetime can give access to many different nodes. The dynamic
relation between iterator values and the nodes they give access to is
exemplified in the extreme by the sorting methods that our lists will provide:
since these operate entirely by modifying links rather than by moving the
contents of nodes, they do not invalidate any iterator value, but the node an
iterator gives access to afterwards is unrelated to the one it gave access to
before.

It is important in practice that in many cases iterators into singly linked
lists are not invalidated by insertion or removal operations, as this allows
several such operations to be done in sequence without requiring a costly
recomputation of iterators to the location in question. Whether the
\Cpp\ standard entertains the possibility that a surviving iterator gives access
to a different element after the operation that before is not entirely clear.
The fact that \std{vector}::\textit{erase} is said to invalidate all iterators
at or beyond the point of erasure [23.3.6.5:3] suggests that the possibility is
not looked upon favourably; this is all the more remarkable because this
stipulation complicates trying to subsequently insert elements at the point of
erasure, even though using the same (officially invalidated) iterator for this
should not cause any problems in practice. On the other hand one would assume
for reverse iterators (say for \std{list}) that they are only invalidated if and
only if their underlying ordinary iterators is; then for them it may well happen
that an operation that does not invalidate them nevertheless makes them give
access to a different item. Indeed our singly linked list iterators behave
exactly like reverse iterators do: they correspond to the separation between two
successive items in the container, and upon structure changes the ``stick to''
the one before (in the sense of the iterator) that separation. In any case,
since the standard seems to allow for deviations from a usual state of affairs
when an implementation technique calls for it (like the number of valid iterator
values for a \frwl.), as long as this does not perturb the proper behaviour of
general algorithms, we fell that the somewhat unusual survival rules of our
iterators under structure modifying operations should not disqualify them.


\subsection{Implementation of simply linked lists with iterators}

We now come to the details of our implementation of simply linked lists and
their iterators. We shall first present it in a simplified form, ignoring for
instance some details needed to make it an allocator aware container, and some
extensions to be discussed later. If one compares with the implementation the
\Cpp\ standard library that comes with our \texttt{gcc} based system, our
implementation is quite straightforward, and close to how we normally write
classes. We found no need for auxiliary (base or implementation) classes, nor
for hidden members with reserved names. Moreover, by implementing links as an
instance of \uptr., memory management is almost entirely handled by the smart
pointer.

First, here is an overview of the class templates implementing the main kind of
simply linked lists (some more definitions that proved useful to include as well
will be mentioned later). There is a node class \slnd., templated over a
type~$T$ for the node contents and an allocator type; it is basically just a
\textbf{struct} containing a link field \textit{next} and a \textit{contents}
field of type $T$. The link has type \lnkt., which is a \textbf{typedef} for
$\uptr.\ang{\slnd.}$ (for now we gloss over the fact that a deleter class is
also specified as second template argument to \uptr.). A
\texttt{const\_iterator} and (derived from it) an \texttt{iterator} class are
defined, both of which have a unique data member $\lnkt.{*}~\lnlc.$. As one
would expect, dereferencing an iterator returns the body
$({*}\lnlc.)\to\textit{contents}$ of the following node, and incrementing
assigns the pointer address $\&({*}\lnlc.)\to\textit{next}$ to the \lnlc. field.
Finally the actual container class template \slls. is defined as a class with
three private data members: a $\lnkt.~\textit{head}$, a
$\lnkt.{*}~\textit{tail}$, and an unsigned integer \textit{node\_count}. The
latter two fields permit efficient implementation of the useful methods
\textit{push\_back} end \textit{size}, respectively.

The user of lists is not given access to entire nodes or any kind of pointers to
them, so the \lnlc. field of iterators is private, but we let all members of
\slnd. be public. Of course the reference to the node body obtained by
dereferencing a \texttt{const\_iterator} is const-qualified, while the reference
so obtained from \texttt{iterator} is not; to enable the definitions of the
necessary overriding of the `$*$' and `$\to$' operators,
\texttt{const\_iterator} declares (its derived class) \texttt{iterator} as
\textbf{friend}. Both kinds of iterator have constructors taking a \lnkt.
argument by reference; it is mainly for internal use by the \slls. class (which
class incidentally is also declared \textbf{friend} of \texttt{const\_iterator},
as its methods must be able to access the links corresponding to iterators
passed to them). The \lnlc. pointer stored in iterators is not const-qualified
at any level: \lnkt. does not have such a qualification, and the access to the
\lnkt. variable itself must be non-const in order to implement link-switching
operations like \textit{insert}. Yet the argument of the
\texttt{const\_iterator} constructor is specified as a reference to
const-qualified \lnkt., so as to make it possible without trickery to form a
\texttt{const\_iterator} in a \textbf{const} method of \slls. like
\textit{cbegin}. Thus the \texttt{const\_iterator} constructor needs to
const-cast the pointer to its argument when setting its \lnlc. field; this is
the only instance of type travesty that our implementation uses.

\subsection{Methods of \textbf{\textit{sl\_list}}$\langle T\rangle$}

The development of our singly linked list class templates was largely demand
driven, starting from a version that had the most basic and common
methods allowing these lists to replace \std{vector} in situation where that
seemed more natural, and adding functionality whenever this served an immediate
purpose and possibly future similar needs. We did have a goal in mind of
eventually having a class template that would provide everything that one would
expect of such templates in the \Cpp\ containers library, and also equivalents
for all of the methods specific to linked lists that \frwl. provides. In the
state we have now arrived at, we feel that this has been achieved; at least for
the kind of use we make of the containers library, our linked lists are on a par
with the containers found there. In addition they have a some specific
functionality that one might not expect in the library; being able to include
such features is one advantage of writing one's own container class templates.

The public methods that \slls. provides can be roughly classified into those one
would expect for a library container class for linked lists (which includes
equivalents for anything the standard specifies for \frwl.), those that extend
these methods but are still within the spirit of the library, and those that add
flexibility to these containers that we had a need for, but which flexibility is
provided in a way not typically found in the standard container library. In the
first category there are some technical methods like \textit{get\_allocator}
related to matters we shall discuss later, which we shall ignore here. We find
in this category the usual panoply of constructors (default, copy, move, from an
iterator-delimited range, from a single element possibly repeated, from an
initializer list), copy-assignment, move-assignment, and the methods
\textit{swap}, and \textit{assign} (which builds a fresh list from its arguments
as constructors do, but assigns that list to an existing variable, overriding
its previous value). Element access, after using the method \textit{empty} to
test if there are any, is provided directly only to the first list item (if
present) by the method \textit{front}, and otherwise requires iterators for
which initial values can be obtained from the methods \textit{begin},
\textit{end}, \textit{cbegin}, and \textit{cend}. One also has modifiers
\textit{clear}, \textit{push\_front}, \textit{emplace\_front},
\textit{pop\_front}, \textit{emplace}, \textit{insert}, \textit{erase},
\textit{remove} (find and remove entries equal to some value),
\textit{remove\_if} (the same with a predicate), \textit{unique} (remove
duplicates), the methods \textit{splice} \textit{merge}, \textit{sort}, and
\textit{reverse} that reorganise lists by manipulating only links, and the
Procrustean method \textit{resize}. Finally the method \textit{max\_size}, of
even more doubtful utility than \textit{resize}, is only provided because \frwl.
does.

So far we have listed only methods that are direct equivalents of \frwl.
methods. The interface to the user of these methods present some differences
with those of \frwl., the most obvious one being the uniform interpretation of
iterator values as designating a boundary between list elements, which is why
the methods \textit{emplace}, \textit{insert}, \textit{erase}, and
\textit{splice} have no \textit{\_\!after} suffix to their name. Also, contrary
to \spla., all our methods called \textit{splice} return an
iterator value indicating the end of the spliced range, which a caller would
otherwise not always have easy access to. This makes is possible for instance to
successively (advancing from front to back) splice in several sub-lists at the
same point in our list; this compares favourably to the effort required to do
the same using \frwl.::\spla.

There are several other methods which, although not provided by \frwl., one
would expect for our list type given that the list object itself has members
informing it of the number of its elements (\textit{node\_count}) and the
location of the final link in the list (\textit{tail}). These methods are
\textit{size}, \textit{emplace\_back}, and \textit{push\_back} (but we do not
provide \textit{back} or \textit{pop\_back}, since there is no efficient way to
do that). In fact the methods \textit{end} and \textit{cend} mentioned before
also use the \textit{tail} member; note that this iterator is not just a marker
for the end of the list as it is for \frwl., but directly allows inserting
elements at the back of the list.

In the category of additional functionality that remains in the spirit of the
containers library, we provide relative versions of the methods \textit{merge},
\textit{sort}, and \textit{reverse}, where the action is limited to a range
delimited by two given iterators. These versions of \textit{merge} (there is one
using \std{less}$\ang{T}$ to compare elements of type~$T$ and another using a
provided comparator) and of \textit{reverse} return the ending iterator of the
merged or reversed range, since this iterator value but would be hard to get at
afterwards if it were not returned. For the relative versions of~\textit{sort},
we deemed it so unlikely that the caller would have any use for the original
value of the ending pointer of the sorted range (which has stuck to the node
that came before it, thereby becoming an iterator somewhere in the sorted range)
that we made this method take its second iterator argument by non constant
reference, and before returning assign the new ending iterator for the sorted
range to it.

We added several methods that use \textit{splice} internally, thus avoiding a
direct call to this somewhat technical method. The methods \textit{append}
(extend at the end by another list), \textit{prepend} (the same at the front),
and \textit{insert} (at a given iterator value) do this, taking the list
argument by rvalue reference so that its elements can incorporated into the
destination list without copying anything. There are also versions of
\textit{append} and \textit{prepend} taking a range delimited by a pair of
iterators (of any type) or an initializer list, like those provided already
for~\textit{insert}. Finally we found that the test for having length~$1$ was
needed sufficiently often to warrant providing a predicate method
\textit{singleton} for it, similar to \textit{empty} for length~$0$.

In the category of additional functionality that one would now expect in a
container library type, let us first mention the predicate \atend. for
iterators, which provides a simple way to test whether an iterator is at the end
of its container (unlike for \frwl., such past-the-end iterators for different
lists do not compare equal among each other, so using iterator comparison
instead, one would need to ensure comparison is with the \textit{end}() of the
appropriate list). We also provide a \textbf{static} method
\slls.$\ang{T}$::\atend. with an iterator argument, which simply calls \atend.
for the iterator; this is done to allow mentioning the list one is expecting to
hit the end of (by calling the method as in
\textbf{if}~(\textit{my\_\!list}.\atend.$(p)$)\dots), but has documentation
implications only. Then we provide \textit{move\_\!assign} and
\textit{move\_\!insert}, each with two iterator arguments, that while traversing
the indicated range will move the values found there into the new nodes rather
than copying them. And since we often use linked lists temporarily to compute
the initial (possibly definitive) value of a \std{vector} variable, we provide
methods \textit{to\_\!vector} that perform the conversion, one of which actually
moves rather than copies the elements into the vector.


\subsection{Flexibility: simpler lists and weaker iterators}

Finally we provide some methods related to variations on lists and iterators,
which we will now discuss. The first variation is a class template \smpl. that
is actually closer to \frwl. than the template \slls. discussed so far.

The class template \slls. and its iterators were set up to provide a singly
linked list container classes with the most general functionality. The fields
\textit{tail} and \textit{node\_count} fields were needed to do that
efficiently, but since they can be recovered from the list structure alone, the
are redundant if efficiency is ignored. Also maintaining these fields does cost
a but if time in many methods, so in use cases where one ends up not using the
methods that depend on these fields, one can gain both time and space by
employing a structure without them. The class template \smpl. provides such a
structure.

The class \smpl.$\ang{T}$ uses the same nodes and the same iterators as
\slls.$\ang{T}$, but has \textit{head} as unique data member of the class
itself. Among the methods of \slls., it implements those that can be done so
efficiently; this means all methods except \textit{size}, \textit{end},
\textit{push\_back}, \textit{emplace\_back}, and all forms of \textit{append}.
We chose to define the method \smpl.::\textit{end} as deleted, rather than to
define it inefficiently by traversing the whole list, so as not to tempt users
to write a traditional looping construct as
$$
\let\&=\textbf \let\\=\textit
  \&{for}~(\&{auto}~\\{it}=\\{list}.\\{begin}(\,);~
     \\{it}\neq\\{list}.\\{end}(\,);~\PP\\{it})~\ldots
$$
which would be quite inefficient (note that \textit{end} is called at each
iteration). Instead one can write
$$
\let\&=\textbf \let\\=\textit
  \&{for}~(\&{auto}~\\{it}=\\{list}.\\{begin}(\,);~
     \&{not}~\\{list}.\\{at\_end}(\\{it});~\PP\\{it})~\ldots
$$
which does the iteration efficiently. As a consequence of the above, whereas
range-based \textbf{for} statement works fine with \slls., it should not be used
with \smpl. since it always essentially translates into it first form
given above.\footnote
{In fact such usage will fail to compile, die to the implicit call of the deleted
 member \textit{end}.}

The methods of \smpl. differ from the corresponding methods for \slls. by
leaving out the parts that deal with \textit{tail} and \textit{node\_count},
thus reducing their work to the bare minimum. It may therefore happen that for a
given list one would prefer for some part of its lifetime that it be packaged as
a \smpl., while for other parts it could be necessary (or preferable) to package
it as a \slls.. Therefore \slls. provides two methods that convert between the
two ways of packaging a linked list, without changing the underlying list
itself: thinking of \smpl. as a naked \slls., we call them dressing and
undressing operations; obviously they transfer ownership of the nodes. Dressing
is done by a \slls. constructor taking a \smpl. argument by rvalue reference
(pilfering it), and undressing by an \textit{undress} method of \slls. that
returns its \smpl. equivalent, while pilfering the list for which it is called.
Since pilfering here means reusing not just value in the nodes but the list
structure itself, it leaves the pilfered list empty. In case one just needs to
know the length or access the final iterators of a \smpl. at some point, we
provide free functions \textit{length}, \textit{end} and \textit{cend} that when
called will find these values (by traversing the entire list), which may be
preferable to converting the list to a \slls. and back again.

We found it occasionally useful to be able to strip down the access to a list
even further and access it through a raw pointer to \slnd., notably for the
purpose of storing it in an object of union type, whose members we prefer to all
be POD values. Such a pointer~$p$ owns the whole list, and should~$p$ disappear
or get overwritten, one should call \textbf{delete}~$p$ first, which as we shall
see takes care of cleaning up the entire list. (Since it is the responsibility
of the destructor of the union type to locate the active variant and to arrange
for its destruction explicitly, this destruction is actually easier when a raw
pointer is stored than if the equivalent \smpl. were.) To this end \smpl.
provides a method \textit{release} that, similarly to the method of
\std{unique\_\!ptr} of the same name, returns such a pointer while emptying the
list itself; it also provides as a constructor from a raw pointer for the
opposite conversion. This allows a list object to ``go under cover'' as a POD
value, and to be resuscitated at a later time. The raw representation is only
intended for use during periods when the list remains unchanged.

A similar but independent flexibility concern is that, when a list (either
\slls. or \smpl.) is to be just iterated over without changing its structure,
one might prefer to use \frwl. style iterators over the iterators discussed so
far, because with one less level of indirection they have a somewhat more
efficient access to the node contents. Thus we define what we call ``weak
iterators'', which are wrappers around just a pointer to a node. Having
different types of iterators giving access to the same container (possibly with
different possibilities for access) is not unusual, as can be seen for
\textit{\textbf{const\_iterator}} versus \textit{\textbf{iterator}}, for $T*$
versus \std{vector}$\ang{T}$::\textit{\textbf{iterator}}, and for normal versus
reverse iterators. Weak iterators (and the related type of const weak iterators)
are forward iterators, and are usable in any context where just requiring
forward iterators are required, but not in those (like the argument of
\textit{insert}) that require either \slls.$\ang{T}$::\textit{\textbf{iterator}}
or \slls.$\ang{T}$::\textit{\textbf{const\_iterator}}. To obtain weak iterators
and const weak iterators, we supplement the \textit{begin} and \textit{end}
methods with counterparts \textit{wbegin} and \textit{wend}, similarly
\textit{cbegin} and \textit{cend} methods with counterparts \textit{wcbegin} and
\textit{wcend}. These methods are defined both for \slls., but also for \smpl.;
indeed \smpl. provides \textit{wend} and \textit{wcend} methods in spite of not
providing \textit{end} or \textit{cend}, as the former just wrap a null pointer
and can be produced cheaply without first traversing the list. (Unfortunately
there is no way to instruct range-based \textbf{for} loops to use alternatives
for the names \textit{begin} and \textit{end}; if this were possible it would
allow using a such loops over \smpl. objects, using weak iterators.) Weak
iterators can even be used to traverse a list that is represented by a raw
pointer to its first node, without first requiring its conversion to an actual
container, since the weak iterator classes have a constructor from such
pointers.

For maximum orthogonality, we define the \atend. predicate for weak pointers as
well, which for them amounts to comparison with an \textit{end()} iterator. In
cases where both iterators and weak iterators are valid possibilities, rewriting
a loop using iterators to one using weak iterators can be as simple as changing
a usage of \textit{begin} to \textit{wbegin}. Therefore it seems best to code
with ordinary iterators initially, and leave the decision to switch to weak
iterators for efficiency reasons to a later stage of program development.

\subsection{Implementing stacks and queues using linked lists}

Singly linked lists are well suited to implement the stack and queue abstract
data types.

The \Cpp standard library provides ``container adaptors'' for these data types:
class templates that take a container class as (second) template argument, and
have a unique data member of that type (a container wrapped by the adaptor).
They have as set of methods those characterising the abstract data type, which
methods call corresponding methods of the wrapped container (often the
corresponding method has the same or similar name; the main point of using the
adaptor is to hide any other methods the wrapped container type might define).
For a stack, the main methods are \textit{push} (or \textit{emplace}),
\textit{top}, \textit{pop} and the predicate \textit{empty}, while for a queue
they are \textit{push} (or \textit{emplace}), \textit{front}, \textit{pop} and
the predicate \textit{empty}. Besides the slight change in the nomenclature
(\textit{top} versus \textit{front}) for inspecting the element that would be
removed by \textit{pop}, the more important difference is that for a stack this
element is the most recently added one among those present (a \textsc{lifo}
regime), while for a queue it is the least recently added one (a \textsc{fifo}
regime). Therefore \std{stack}$\ang{T,\textit{Container}}$ maps its
\textit{push}, \textit{top}, and \textit{pop} methods respectively to
\textit{push\_back}, \textit{back}, and \textit{pop\_back} methods of
$\textit{Container}$, while \std{queue}$\ang{T,\textit{Container}}$ maps its
\textit{push}, \textit{front}, and \textit{pop} methods respectively to the
methods \textit{push\_back}, \textit{front}, and \textit{pop\_front} of
$\textit{Container}$. The \textit{Container} class must provide the mentioned
methods under these names. When used without specifying a container type as
template argument, these adaptors will be based on the \std{deque}$\ang{T}$
container by default.

Given these requirements, it is possible to use \slls. as underlying container
type for a queue, forming the type $\std{queue}\ang{T,\slls.\ang{T}}$, since
\slls. defines methods \textit{push\_back} (and \textit{emplace\_back}),
\textit{front}, \textit{pop\_front}, and \textit{empty}; on the other hand, the
class \smpl. lacks a method \textit{push\_back}, so it cannot be used to
implement a queue. For \std{stack} on the other hand, there is a difficulty in
that it wants to use \textit{back} and \textit{pop\_back} methods of the
container (a choice probably motivated by the idea of using \std{vector} as
container), which neither of our singly linked list containers define. But
singly linked lists (even of the \smpl. kind) are very well capable of
implementing stacks: they just need to add and remove nodes at the front of the
list rather than at the back. Our remedy for stacks is then to define a class
template \mspl., which is an adaptation of \smpl. designed to be used
(exclusively) as template argument to \std{stack}. It is publicly derived from
\smpl., and apart from its constructors, it only defines methods
\textit{push\_back}, \textit{emplace\_back}, \textit{back}, and
\textit{pop\_back} which respectively forward to the methods
\textit{push\_front}, \textit{emplace\_front}, \textit{front}, and
\textit{pop\_front} of its base class. (As it defines only the methods that
\std{stack} uses, but inheriting everything else from~\smpl., the type \smpl. is
a rather curious container type, for which for instance \textit{front} is
synonym for \textit{back}. It being derived from \slls. is only relevant for
classes derived from $\std{stack}\langle T,\mspl.\ang T\rangle$: they can
directly access the container object, which is a \textbf{protected} member of
\std{stack}.) For completeness we also define a type \msll. similarly derived
from~\slls..

There is a slight difficulty with using an instance of \slls. as container type
for \std{queue}, since this adaptor not only defines the main methods of a queue
abstract data type listed above, but also \textit{back}, which it forwards to
the method of the same name of the container; but \slls. has no such method, nor
can it be efficiently implemented without recording additional data. However, we
consider this to be an error on the part of the \Cpp\ specification, which
should not have defined a \textit{back} method for \std{queue}, since access to
the most recently pushed item is not part of the queue abstract data type. Indeed
in none of the many places where we use queues in our software do we have any
use for such a method. We have chosen to use $\std{queue}\langle T,\slls.\ang
T\rangle$ in these places; since methods of a class template that are never
called are not compiled either, this causes no problems. Similarly
$\std{stack}\langle T,\mspl.\ang T\rangle$ cannot handle calls to \textit{size},
but code using a stack should not have any need for such calls, and in the
absence of such calls the program will compile correctly.

If it were necessary to have a container type that can be used with \std{queue}
according to the full specification of the \Cpp\ standard, one could easily add
a data member to \mspl. that allows tracking the most recently added item in all
its methods; using that implement \std{queue}::\textit{back} efficiently when
using this container. But it seems contrary to \Cpp\ design principles to burden
users with extra space and time cost in order to support functionality that few
will ever need. On the other hand there is some functionality compatible with
the abstract data type that one can add at no cost to that of \std{queue} and
\std{stack} when using a simply linked list as underlying container: a
constructor that takes a \std{initializer\_list} (often a singleton), and an
alternative to the \textit{pop} method that splices the element popped directly
into a \slls. or a \smpl. (avoiding the cost of constructing a new node, and the
need to call \textit{front} or \textit{top} before doing~\textit{pop}).
Therefore, rather than overriding the default underlying container every time a
queue or stack variable is introduced, we provide (in our namespace) templates
for classes \queue. and \stack. that are \emph{derived} from
\std{queue}$\ang{T,\slls.\ang{T}}$ respectively from
\std{stack}$\ang{T,\mspl.\ang{T}}$, and which provide these extra methods that
we found useful. The also override as \textbf{delete} the methods \textit{back}
and \textit{size}, respectively, which they inherit but would forward to
nonexistent methods of the container type used, so that any attempt to use them
will elicit a clear error message.

\section{Practice}

The functionality for singly linked lists described above, was developed and
implemented over the course of time, in the context of the project Atlas of Lie
groups and Representations (which we shall henceforth abbreviate to Atlas); the
commit where the class templates for singly linked lists were first introduced
is dated June~20, 2014. Specification and implementation were gradually refined
and enriched over time. At some point a clear and prolonged effort was made
though to completely cover the functionality of \frwl., retaining proximity to
its specification, and to conform as much as possible to the standard library
conventions for container types. In this section we shall discuss how the
container type serves the needs we have encountered in withing our project, and
how its utility compares to alternative ways available to meet those needs.

\subsection{The Atlas project}

Let us start describing the goals of the Atlas project. It is a project in pure
mathematics, which aims at developing understanding of the representation theory
of real Lie groups, with the aid of software that implements the various
relevant algorithms, mostly of combinatorial nature, in a sufficiently efficient
manner to be applied to very substantial problems. Different types of Lie
groups, different representations, and different questions asked about those
representations provide a wide variety of computational challenges. The ultimate
goal that was envisioned when the project was founded in 2002 is finding a
complete description of the unitary dual of the split group of type~$E_8$. While
this goal has not been achieve currently, substantial progress has been made
over the years, including complete results for necessary ingredients for the
$E_8$ computation (notably the determination in 2007 of the full table of
Kazhdan-Lusztig-Vogan polynomials), and the full unitary dual
description for somewhat smaller groups; as computer power and storage
capabilities have substantially increased over the years, it would appear that
the mentioned ultimate goal should be within reach in the foreseeable future.

To perform the computation, the project develops a special purpose computer
algebra package, originally conceived and implemented by Fokko du Cloux, and
further developed after his premature death under the technical direction of the
current author. The software executable has two major components of rather
different nature: on one hand an extensive library of specialised algorithms
implemented in \Cpp, and on the other hand a user programming interface that
allows algorithm development by mathematically oriented users. A quite extensive
collection of programs written in this user programming language has been
developed in addition to the \Cpp\ executable.

The \Cpp\ library centres around a number of discrete, graph-like structures
(root systems, double cosets $K\backslash G/B$, blocks of representations) that
are generated dynamically, and whose vertices with associated data are stored in
tables (implemented mostly using \std{vector} instances), the mathematical
values represented by such a vertex being henceforth most often represented by
just an index into that vector. The user programming interface is a compiler and
interpreter for a language somewhat resembling Python (but with static typing
and not ostensibly object oriented). It has a layered structure common to many
language processing tools: input and output facilities, a lexical analyser, a
(generated) parser, a type checker that also transforms expressions and
functions into executable tree structures, and an evaluator for such structures.

The implementation of both parts was firmly in place long before we introduced
our linked list container class template, and so the pieces of code where such
lists are now used were either added recently, or (more often) have been adapted
to use a list where previously some other solution was used. That other solution
was often to use \std{vector} instead, in the case of stacks and queues to
(implicitly) use \std{deque}, and in some cases to use specialised recursive
definitions of linked structures. (While there were certainly places where using
\std{list} would have been an option, it was in fact never done. As for \frwl.,
although it had become officially part of the language by the time we introduced
\slls., we would at that point still deemed using it a portability liability,
so it was never really considered.)

\subsection{Use of singly linked lists in the algebra library}

We only introduced \slls. (or sometimes \smpl.) in places where this was more
convenient to the programmer or led to a simplified and no less efficient
implementation. This includes all uses of stacks and queues, for which we used
the class templates we based on singly linked list described above. This was a
drop-in replacement, with a small added convenience that we added for placing
one or more values into the data structure at construction. The change leads to
a substantially more modest memory footprint, especially in cases like ours
where the size of the stack or queue typically remains small (a \std{deque}
needs around $600$ bytes of storage even when it is empty, as opposed to about
$20$ bytes for an empty \slls. instance) and simpler implementation of
operations. In the Atlas library, these data types are used in about a dozen
places, generally to hold yet-to-process elements in some orbit generation
process or similar kind of closure operation. Either a stack or a queue would do
for this purpose, with a stack giving more of a depth-first approach versus a
queue leading to more breadth-first processing; we generally prefer the latter,
as we expect it to give a smoother evolution of the storage needed.

Apart from these cases, there are a number of places in the Atlas library where
we decided to use lists instead of vectors when a sequence of values needs to
be computed, and its length cannot be predicted beforehand. Whether or not lists
are better or worse from a performance point of view is not the issue here: the
old solution worked perfectly well, and since these are places where the program
spends a quite insignificant fraction of its time, any effect that there might
be would be impossible to measure. But the simple and straightforward task
matches the characteristics of linked lists perfectly, while using vectors one
knows that some effort (copying entries while growing) and space (in the final
result) are being spent that are not strictly required for that task; the list
solution avoids for the programmer distracting thoughts about some possible
minute improvement.

In a few cases, using a list rather than a vector to accumulate a sequence of
values allowed code to more fluently express the operation being implemented. In
particular we mention a normalisation process, in which for a certain class of
values, some may be rewritten as combination of $0$, $1$, or $2$ other values,
until being left with only ``final'' values that cannot be rewritten (which is
ultimately bound to happen). The original implementation handled this in a
somewhat clumsy manner, trying to produce a vector of values in a single sweep
(repeatedly calling \textit{push\_back}), but having to revert to appending one
vector (obtained from a recursive call) to another one in certain cases. Using
lists we can follow the same logic, but the appending operation is now painless:
the \slls.::\textit{append} method is efficiently implemented by just
redirecting a link. But using linked lists in fact obviates any need to worry
about finding a good organisation of the rewriting process in the first place:
just about any approach can be realised efficiently, including ones in which
values are being rewritten while sitting in the middle of a list that is being
processed. So linked lists meets the requirements of the task particularly well
here, although again this is more about elegance than efficiency, since the
rewriting process is unlikely to lead to a large set of values anyway. So
in conclusion, having a singly linked list container type available in the Atlas
library occasionally allows a more satisfactory solution than before, even
though the improvements are by no means a spectacular.

%Exporting a list-valued result of variable (and unpredictable) length
%structure/dynkin.h:  containers::sl_list<RankFlags> components() const;
%structure/innerclass.h:containers::sl_list<TorusPart> preimage
%structure/innerclass.h:containers::sl_list<TorusPart> central_fiber
%gkmod/blocks.h:  containers::sl_list<BlockElt> finals_for(BlockElt)
%gkmod/repr.h:  containers::sl_list<StandardRepr> finals_for(StandardRepr)
%gkmod/ext_block.h:containers::sl_list<std::pair<..,..> > extended_finalise
%gkmod/ext_block.h:containers::simple_list<BlockElt> condense(Matrix&,..)
%gkmod/ext_block.h:DescValue star(...,containers::sl_list<param>& links);
%gkmod/repr.cpp:  containers::simple_list<BlockElt> finals = ...
%gkmod/ext_block.cpp: containers::sl_list<param> links; // receive |star| output
%gkmod/ext_block.cpp: { containers::sl_list<param> links; // same
%gkmod/ext_block.cpp:  containers::sl_list<param> links; // idem


%Special case: normalising a term through certain {0,1,2}-valued descents
%gkmod/repr.cpp:  containers::sl_list<StandardRepr> result { z };


%Use as stack or queue while generating an orbit of equivalence class
%structure/realredgp.cpp:{ containers::stack<std::pair<TorusPart,Grading> > to_do;
%structure/weyl.cpp:  containers::queue<WeylElt> to_do { w };
%structure/weyl.cpp:  containers::queue<TwistedInvolution> to_do { tw };
%gkmod/standardrepk.cpp:  containers::queue<KGBElt> queue { root };
%gkmod/standardrepk.cpp:  containers::queue<equation> queue;
%gkmod/standardrepk.cpp:  containers::queue<q_equation> queue;
%gkmod/blocks.cpp:  containers::queue<BlockElt> queue { size() }; // invol. packet
%gkmod/blocks.cpp:   { auto to_do = containers::queue<KGBElt> { ctx1 };
%gkmod/standardrepk.cpp:  containers::queue<KGBElt> queue { root };
%gkmod/standardrepk.cpp:  containers::queue<equation> queue;
%gkmod/standardrepk.cpp:  containers::queue<q_equation> queue;
%gkmod/ext_block.cpp:  containers::queue<param> to_do { param(ctxt,sr) };

%Intermediate storage when accumulating a vector
%structure/rootdata.cpp:  containers::sl_list<weyl::Generator> w;
%structure/rootdata.cpp:  containers::sl_list<weyl::Generator> w;
%structure/rootdata.cpp:  containers::simple_list<weyl::Generator> w;

\subsection{The interpreter}

While the user interface of the Atlas software originally just consisted of code
to gather input for library functions and to print their results, this has since
been replaced by a programmable interface, which over time has become a
full-fledged programming language, representing a substantial part of the code
base. In the interpreter, linked structures are much more frequently used than
in the library part, so there was much more opportunity for using linked lists
in this part than in the rest of the software.

This does not apply to the first stages of input processing, namely gathering
lines of input and lexical scanning. Here the only point that might be worth
mentioning is that the possibility to include at some point of a file the text
from another file leads to a stack of active input sources being maintained,
where we use the technique mentioned above of defining a derived
not-quite-abstract data type of a stack enriched with extra methods. The extra
methods allow finding the current size, as well as obtaining constant-iterators
for traversal (which is needed to implement the safety precaution that no file
can be opened a second time while it is already an active input source); as
underlying container for this stack we use \textit{mirrored\_sl\_list}.

Our parser is produced using the \texttt{bison} parser generator program. Thus
we have no control over the parser function itself, which is a table-driven
skeleton program written in \Cee, but we do provide its action functions which
build up parse tree of the user input as it is being recognised. The generated
parser can actually by compiled as \Cpp\ code, so the action code is not
restricted to what can be written in \Cee, and notably the parse tree can use
\Cpp\ types. These types have mutually recursive definitions that more or less
reflect the context-free grammar being implemented, using pointers to refer to
each other, and using a tagged \textbf{union} of variants when a type represents
a syntactic category with many different production rules.

\section{Implementation details}

Implementing various methods for manipulating singly linked lists is an
instructive exercise, resulting in most cases in quite straightforward code.
Yet there are a few occasions where our implementation seeks a slight
improvement to the most obvious approach;  we briefly mention here a few of such
cases.

\subsection{Details for methods of \textbf{\textit{sl\_list}}$\langle T\rangle$}


The implementation of the container class template \slls. itself, though long,
is quite simple; in particular it ensures exception safe memory management
without needing any contortions. As we mentioned, \lnkt. is a \uptr. instance,
and this takes care of almost all such concerns, and it does so with perfect
timing. For instance, the destructor of \slls. has empty body (it can be
defaulted) since destruction of the \textit{head} field will release memory for
the first node if there is one, and as this begins with destructing the node,
will also recursively clean up the rest of the list. (We shall mention below a
small adaptation to avoid the deep implicit recursion this would imply for long
lists, but barring possible run time stack overflow, things would work just fine
without the tweak.)

To erase a node, given an iterator \textit{pos} that currently gives access to
it, we make an alias of the link it points to as
$\lnkt.\&~\textit{link}={*}\textit{pos}.\lnlc.$, and then redirect that link by
calling $\textit{link}.\textit{reset}
(\textit{link}\to\textit{next}.\textit{release}())$. Here is why this works. The
call of \textit{release} fetches the link in the node to be erased, returning it
as a raw pointer~$p$ while replacing that link itself by a null pointer; then
the \textit{reset} method replaces \textit{link} (after recording its old value)
by a new smart pointer reconstructed from~$p$, and finally deletes the old value
of \textit{link}, which calls the destructor for the node. Since at the time of
its destruction the linked-out node contains a null pointer in its \textit{next}
field, this one node that becomes unreachable is destroyed, but no further nodes
are. The same is true when one links out a whole range of nodes at once; as
illustration here is the complete code of the method \slls.::\textit{erase} that
does this.
\begin{verbatim}
  iterator erase (const_iterator first, const_iterator last)
  {
    node_count -= std::distance(first,last);
    first.link_loc->reset(last.link_loc->release()); // link out range
    if (at_end(first)) // whether we had |last==end()| initially
      tail = first.link_loc; // reestablish validity of |tail|
    return iterator(*first.link_loc); // convert |first| to |iterator|
  }
\end{verbatim}
Note that although it might seem that the call to \std{distance} in the first
line is the only part that takes time linear in the length of the range, this is
also the case for the second line linking out the range: it implicitly destroys
each node in the range. That second line could incidentally also have been
written using a move-assignment, namely as $*\textit{first}.\textit{link\_loc}
=\textit{std::move}(*\textit{last}.\textit{link\_loc})$. It might not be obvious
that the move assignment still ensures correct timing, notably that
$*\textit{last}.\textit{link\_loc}$ is null when the old
$*\textit{first}.\textit{link\_loc}$ is destructed, but in reality the semantics
of this move-assignment is given by precisely the combination of
\textit{release} and \textit{reset} calls above, so the alternative here is just
one of form, not of substance.

Inserting a node into a list requires a few more operations, but the most
straightforward approach is still correct: make a local \lnkt. variable point to
a newly allocated node, move the link to what should become its successor into
its \textit{next} field, and finally move the link from the local variable to
the location where it belongs. This is exception-safe: the only operation that
could throw an exception is the initial allocation, and should this happen, the
list is still in its initial state, so it either survives, or it is cleaned up
without leaks or other anomalies (depending on where the exception is caught).

***

Few of these methods present any difficulties or surprises in their
implementation. A small point of attention is that move construction or
assignment should not simply move all fields, but in case of an empty list make
\textit{tail} point to the \textit{head} link in the own list, not the list
moved from. Another point to pay attention to is that when a node (or range) in
front of a given iterator is removed, that iterator remains valid and henceforth
accesses the node after the removal, if any; therefore methods like
\textit{remove} that traverse a list while removing some nodes \emph{do not}
increment their iterator of traversal in the case where a node is removed. This
is similar to what \frwl. requires, but contrasts with the situation for
\std{list} where the iterator passed to \textit{erase} is invalidated by the
call, whence a call of the form $\textit{erase}(p\PP)$ is needed. The latter
idiom has no use with our \slls. (it would just make $p$ invalid after the
erasure), and to avoid its accidental use we prefer to not declare any
post-increment method for iterators at all.

Maintaining \textit{tail} correctly in all methods requires
some care, but never requires more than constant time. Maintaining
\textit{node\_count} can usually be done by keeping track of any nodes that are
added or removed, but on two occasions it involves computing the length of an
entire range (which takes time linear in that length): one is that cited above
of erasing a range, and the other is when using \textit{splice} to transfer a
range from \emph{another} list into the current list. While the former operation
would take linear time anyway, this is not the case for \textit{splice}: to
modify the list it suffices to cyclically permute three links, which takes
constant time. So this case, exceptionally, the complexity of the method is
dominated by the need to maintain \textit{node\_count}. We note in passing that
due to the way the $3$-iterator instance of \frwl.::\textit{splice\_after} is
specified, it also has complexity linear in the length of the spliced segment.

Special caution is required for methods that operate on a range inside a singly
linked list and which (potentially) change its structure; this concerns the
``range'' versions of \textit{reverse}, \textit{merge}, and \textit{sort} (all
of which are overloads that we added, not present for \frwl. or \std{list}).
After the calling such a method, the restructured range may no longer be
delimited by the iterators passed to the method. Since iterators stick to the
node before them, the starting iterator of the range still marks the start of
the restructured range, but the ending iterator is no longer necessarily at its
end (it need not access the node after the range, or nothing if that was the
case before restructuring). For this reason we let these methods return the
ending iterator of the restructured range, so that the caller may take note of
the change, or in the case of \textit{sort} (where we cannot imagine any caller
not needing to record the bounds of the sorted range) we even let the method
take the ending iterator by reference, and assign the new ending iterator to it.


\subsection{Destruction of lists}

As was already mentioned, destructing a link (by calling \textit{reset} or by
assigning to it) will destroy the \textit{sl\_node} it points to, and
recursively all other nodes reachable from it. In general this is precisely what
is desired, but letting this recursion proceed naturally would lead to nesting
of active destructor calls to a depth given by the length of the list being
cleaned up, which for sufficiently long lists might cause the \Cpp\ runtime
stack to overflow and the program to crash. To avoid this, we define the
destructor of \textit{sl\_node} so as to achieve the same effect, but using a
loop rather than through recursive calls. In the destructor body we repeatedly
remove the following node as long as there is one, each time calling
$\textit{next}.\textit{reset}(\textit{next}\to\textit{next}.\textit{release}())$
similarly to what \textit{erase} does; then we let our destructor complete,
which will finish by calling the destructors for the fields of the
\textit{sl\_node}, but \textit{next} by now assured is to hold a null pointer
so that its destruction is trivial. Each of those initial calls to
\textit{reset} does lead to a recursive call of the \textit{sl\_node}
destructor, but for a node whose \textit{next} field cleared so that recursion
is limited in depth to two levels.

\subsection{Usage of structure modifying methods}

The iterator passed to \textit{insert} will remain valid after
the insertion, and will then refer to the beginning of the inserted range. It is
then natural that the iterator returned by the call should be the one
corresponding to the end of the inserted range; then calling the same method
again with the returned iterator causes an insertion after rather than before
the first insertion. This does not satisfy the sequence container requirements
for \textit{insert}, which want the returned iterator to refer to the
\emph{start} of the inserted range, but is similar to the specification of
\insa.. The return value of \textit{erase} is less relevant, as in all cases it
equals an iterator passed to the method, but the natural iterator to return is
the one indicating where nodes were erased (the place where nodes should be
inserted if the removal were to be undone). That agrees with sequence container
requirements (by contrast note that the iterator returned by \ersa. is not the
one that should be passed to \textit{insert\_after} to undo its effect), and
means in our case that it is a copy of the starting iterator for the deleted
range. Finally, it is similarly natural that \textit{splice} should return an
iterator corresponding (as fence post) to the end of the inserted sequence; that
is usually the ending iterator of the range passed to \textit{splice}, but when
the range happens to be empty it is the iterator passed as the point to splice
into. (Curiously \frwl.::\spla. returns nothing at all; this makes it harder to
locate the end of the spliced range, for instance to splice a second range after
it. Maybe this is because specifying the correct iterator to return, namely the
one for which a second \spla. would continue after the first range, would
require some awkward wording in the standard to cover all cases correctly.)

\subsection{Possibilities to avoid node allocation and deallocation}

When the \textit{assign} method must replace the contents of a list with a new
one (for instance with a copy of a range of some other container), one could
just destroy the old list and generate the new list as upon construction;
instead we prefer to overwrite the contents of existing nodes as much as
possible, and then if the lengths do not match either truncate, or extend with
the remainder through a call of \textit{append}. This \textit{assign} method is
itself called upon by the copy-assignment operator, which therefore does the
same optimisation.

When inserting nodes, there is a slight difference according to whether this
happens at the end of the list or not: one needs to advance \textit{tail} in the
former case only, but on the other hand no work is needed there to attach the
tail of the list to newly created nodes (a null link is set by our \slnd.
constructor, and it can remain in this case). Moreover, when adding nodes at the
end we can use the \textit{tail} field where we would otherwise need a separate
iterator. All told, an insertion at the end requires fewer operations per node
than one at another position. In view of this, we implement \textit{insert} of a
range either by a calling \textit{append} (when the insertion is at the end), or
else by generating a temporary list with just the inserted range, splicing it
afterwards into our list at the proper place.

\subsection{Merging lists}

The \textit{merge} of two sorted lists can be described in a manner symmetric in
the two lists: repeatedly take the smaller of the head nodes remaining in the
lists as long as both exist, and finally add what remains of the list that did
not run out. Instead we use a less symmetric description (viewing it as another
list being merged into ``our'' list) that minimises the number of link
operations to be done: traverse our list one node~$x$ at a time, each time
isolating from the other list the sequence of all remaining nodes that compare
less than~$x$, and if non empty splice that sequence before~$x$. Then in case
the other list completely ran out, we return from \textit{merge} immediately
after splicing the sequence; if our traversal completes without such a return
happening, this means some nodes remain on the other list, which remainder we
then splice to the end of our list before returning.

\subsection{Allocators}

Container class templates are expected to be allocator aware, which means they
take an allocator type as template parameter, and containers have as a component
(most often as a base class) a value of that allocator type; all storage for the
container itself (as opposed to additional storage possibly required by its
elements) should be acquired and released using methods of the allocator. While
the standard lacks any clear and detailed motivation, the idea behind allocators
is probably to provide the user a hook to improve upon the inefficiency that
could result from calling the default allocator, and thereby indirectly the
operator \textbf{new} and possibly the low level \textit{malloc} function, every
time the container class needs to acquire new storage, and similarly for
disposing of storage.

Although the idea seems straightforward, there are constraints and subtleties
that make allocators hard to grasp, and it appears that allocators other than
\std{allocator}, the provided default, are rarely used. People do mention that
custom allocators designed for specific applications may speed them up
considerably, but few widely useful allocators seem to exist. Searching for
information on the subject, we perceived a supposed tension between being
general and being efficient; however, the standard provides little in the way of
tools to build allocators taking advantage of specific situations, like the fact
that many container types always allocate storage in units of fixed size. The
standard does give an extensive list of requirements for allocators to be fit
for use with container classes.

An allocator type $A$, necessarily a \textbf{class}, can only allocate for a
specific type, given by the member type $A::\textit{value\_type}$; it can
however be asked ask for contiguous storage for $n$ objects of that type. But
allocators for use with containers should not be defined to handle just one
value type, since a container for values of type $T$ (a template argument of the
container type) will receive as parameter an allocator for $T$, but often needs
to allocate storage for a different type; in our case of $\slls.\ang T$ for
instance, allocation is for $\slnd.\ang T$ rather than~$T$. This almost forces
allocator types to be defined by class templates, themselves taking
\textit{value\_type} as template parameter. When that is the case, the template
instance with another type~$U$ as \textit{value\_type} is obtained as
$\std{allocator\_traits}\ang A::\textbf{template }\textit{rebind\_alloc}\ang U$,
which is what a container class implementation should use as allocator type when
it needs to allocate for another type than the container's value type. But that
allocation requires an actual value of that new type~$B$; to make that possible
the standard requires that cross-type copy-construction be possible, namely that
a declaration $B~b(a)$ copy-constructs an allocator of type~$B$ from an
allocator~$a$ of type~$A$. This makes it hard for allocators to have any non
static data members with a type depending on their \textit{value\_type}; say for
example that $a$ would hold a pointer to storage it has ready for future
allocations for type~$T$, how should the corresponding member of a
copy-constructed allocator for an unrelated type~$U$ be initialised? Static data
members do not have this conundrum (constructors can ignore them) and are more
suited to store resources of the allocator that are specific for~$T$. We may
conclude that allocators probably are either stateless (having no non static
data members, so that all objects of the same type are equivalent) or manage
their resources in a way ignorant of value types (if any copy constructed object
shares those resources with the original); incidentally, the default allocator
fall into both categories at once.

In any case, even if allocators remain a somewhat elusive species, our container
class template must accommodate their use, and (as of \Cpp11) the possibility
that they have state. But a first concern is making \uptr. work correctly with
allocators. Now this smart pointer class template pursues a somewhat different
abstraction than that of allocators: it manages ownership of some object that
may or may not be a pointer, and since ownership is all about cleaning up when
one is done, it incorporates a ``deleter'' object, of type given by a second
template argument to \uptr., that is called to perform cleaning up. For the
common case that the managed object is in fact a pointer, a template
\std{default\_delete} for a stateless deleter class is provided, whose disposal
action is to call \textbf{delete} for the pointer given to it. But although the
storage given out by an allocator is most likely accessed through a regular
pointer, calling \textbf{delete} for it would bypass the proper way to free the
storage, namely to call the \textit{deallocate} method of the allocator.
Strangely nothing appears to be provided in the standard library for this case
(although a proposal~\cite{P0316R0} was submitted in 2017), so we wrote a small
class template \textit{allocator\_deleter} such that
\uptr.$\ang{T,\textit{allocator\_deleter}\ang{\textit{Alloc}\ang T}}$ can be
used to manage pointers to $T$ obtained from the allocator $\textit{Alloc}\ang
T$. Of course we must also make sure to get a pointer from the allocator to
begin with, so we also wrote a function template \textit{allocator\_new} to
replace \textbf{new} for obtaining such pointers. Some care was needed to make
it safely handle, just like \textbf{new} does, the possibility of an exception
being thrown during the construction of the value in the allocated storage.

The pointer returned by \textit{allocator\_new} should always be used to
initialise an instance of \uptr.. Our code indeed does so (with that instance
being \textit{link\_type}), though not always instantly: in several cases we
move a link into the \textit{next} field of the newly allocated node (which
cannot throw any exception) before constructing a link from the mentioned
pointer, typically by passing it to the \textit{reset} method of the link that
was just moved from. There are other ways to go about this, and many would
consider it more robust to wrap \textit{allocator\_new} in a function returning
a \uptr. right away; that is possible, and which option one prefers depends on
whether one has more confidence in a statement that cannot throw not throwing,
or in the code generator optimising away the part of the \uptr. functionality
that can be seen to never do anything.

With these preparations done, most of the allocator awareness is taken into
account automatically. There just remain some questions about what to do with
contained allocator objects (that might have state) when copying, moving or
swapping containers. These questions could be ignored prior to \Cpp11, as
container classes were allowed to assume that allocators of the same type always
compare equal, which means that storage allocated by one allocator object can be
correctly deallocated by another. In \Cpp11 more legitimacy was accorded to
stateful allocators, which might not support such cross-allocator recycling; if
indeed they do not (the allocators compare unequal), then the container
implementation must take care to ensure that at all times the storage and
allocator in a same container remain compatible, so that safe destruction is
possible.

Also allocators now control several aspects of the questions arising, through
members of their \std{allocator\_traits} instance: how to obtain an
allocator object during copy construction of the container (reasonable choices
are copy-constructing the other allocator or default-constructing), and whether
during copy-assignment, move-assignment and swap of containers, the same
operation should be done with the contained allocator objects. (No choice is
given for move-construction: here the allocator will be moved, period.)

Once one understands what these traits members are about, it is easy to get the
implementation to respect the choices they express, but it does complicate the
code for assignments and swapping. For each case there is an easy scenario, when
one proceeds as if allocators were no issue, and a harder one. For
copy-assignment, the easy scenario is when the old allocator remains in place:
one can overwrite existing node contents as described above for \textit{assign}.
But when the allocator must be propagated, one has to first deallocate the old
container contents (using the old allocator), then copy-assign the allocator,
and finally construct (using the new allocator) new storage to hold a copy of
the list. For move assignment and swap however, the easy case is propagation of
the allocator object: one can just move-assign respectively swap the
allocator(s) and all members of the \slls. object(s), just taking care to adapt
\textit{tail} fields for empty lists. (Doing so, allocators mismatch container
contents for a short time, but no exception can be thrown.) But when
propagation is forbidden and allocators differ, the best one can do then is to
move/swap individual list entries up to the length of the shorter of the two
lists, and copy into freshly allocated nodes for the remainder. The moral is
that life is easiest when allocators follow ownership (transfer in case of move
semantics, but not when copying). As the same would seem to be true for other
container types as well, one wonders why the standard does not just impose the
choices that give rise to easy scenarios; the possibility to copy-construct a
container with explicitly provided allocator would still give the user full
control.


\begin{thebibliography}{Comph}

\newcommand\tit[1]{\newblock ``#1'',}
\newcommand\journ[3]{\newblock \textit{#1}, \textbf{#2} (#3),}
\newcommand\publ[2]{\newblock \textit{#1}, #2,}

  \bibitem[Atlas]{Atlas source}
  \tit{Atlas of Lie groups and Representations} software project,
  \url{https://github.com/jeffreyadams/atlasofliegroups}; the
  implementation of singly linked lists is in subdirectory
  \texttt{sources/utilities} of the repository, in file \texttt{sl\_list.h}

  \bibitem[Comph]{triple ref}
  Computerphile YouTube channel: \tit{Triple Ref Pointers}
  \url{https://www.youtube.com/watch?v=0ZEX_l0DFK0} (2017)

  \bibitem[Knejp]{P0316R0}
    \author{Miro Knejp}
    \tit{\texttt{allocate\_unique} and \texttt{allocator\_delete}}
    ISO/EJC JTC1/SC22/WG21 - Papers 2017, proposal P0316R0
    \url{http://www.open-std.org/jtc1/sc22/wg21/docs/papers/2017/}

  \bibitem[Knuth]{TeX the program}
  \author{D.~E. Knuth}
  \tit{\TeX, the program}

\end{thebibliography}
\end{document}

%Local IspellDict: british

% LocalWords:  dereferenceable
